\documentclass[11pt]{article}
\usepackage[T1]{fontenc}
\usepackage[utf8]{inputenc}
\usepackage[margin=25mm]{geometry}
\usepackage{booktabs,longtable,array,amsmath,url}
\usepackage[hidelinks]{hyperref}
\setlength{\emergencystretch}{2em}
\title{Supplement: Final ContactAwareSS Evidence and Reproduction Contract}
\author{Mukhammed Zhumali}
\date{25 October 2026}
\begin{document}
\maketitle

\section{Executable and Evidence Identity}
This supplement accompanies the separate final manuscript. Prior SE2, R1 and PC2 sources,
results, papers and archives are preserved as historical evidence. They are not relabeled as
new final-campaign observations. All 108 final configurations use the same Java executable.

\begin{longtable}{>{\raggedright\arraybackslash}p{.28\textwidth}>{\raggedright\arraybackslash}p{.66\textwidth}}
\toprule Artifact & Identity \\
\midrule\endhead
Base Git commit & \texttt{3621333d8093b4cb0c232}\newline\texttt{c7db9164c66e273053c}; the working tree was dirty, so this commit alone does not identify the executable. \\
Source manifest & 181 Java/build/wrapper files, SHA-256 \texttt{71c73f5e99550ee6365556fcebcc9411405}\newline\texttt{32dea2cf7a0193de7494bbbde4d07}. \\
Final-v2 JAR & SHA-256 \texttt{bb694a9e1c3341ab7ab9602348a34a934}\newline\texttt{0daf9fcf708d6906435dcaeb35439c6}. \\
Registered protocol & SHA-256 \texttt{2c406eb1c13dd88cea8fc7731736678e8}\newline\texttt{fd3f5fae2b343a48395b5666a878cb8}. \\
Build & Java 21.0.10; Maven 3.9.9; isolated clean verify. 478 tests, 64 test classes, zero failures/errors/skips. \\
\bottomrule
\end{longtable}

The first isolated full-suite attempt failed five HTTP-import tests because the sandbox denied
loopback connections; it had zero assertion failures. A second clean build with local HTTP
permitted passed the full suite. Both build records are retained, and their source manifests
are identical. No Java change or test exclusion was used to obtain the final JAR. Source ZIP
timestamps can differ in a rebuild; the packaged JAR provides the exact binary used here.

The machine records are \texttt{target/final-research/manifest.json},
\texttt{build-validation.json}, the two build logs and Surefire XML. Two compact PC2 oracle
files separately record COMPLETE at 1364 nodes under each tie rule. The older small-quality
test also retains bounded-controller differences; a passing test suite is not a claim that
every bounded policy attains every oracle optimum.

\section{Fixed Matrix and Input Interpretation}
\begin{center}
\begin{tabular}{lr}\toprule Block & Attempts \\
\midrule
X5/Y/prior-demand regression, offline and measured & 18 \\
Three missions, three loads, six offline configurations & 54 \\
Three ISS loads, five measured configurations & 15 \\
INSTALL/cryptoperiod/horizon sensitivity & 9 \\
Disclosure and original temporary channel loss & 10 \\
Initial management-operation interruption & 2 \\
\midrule Total & 108 \\
\bottomrule\end{tabular}
\end{center}

The full chronological schedules contain ISS46, NOAA38 and HST48 contacts. Different missions
use different saved dates; this is not a simultaneous operational mission comparison. Public
GP, station-coordinate and Earth-orientation files supply geometry. Satellite-station
authorization, payload availability, cryptographic operation timing and lifecycle limits are
model assumptions. The schedule selector is independent of security-policy results.

The common rate is 100~B/s with 1000~B complete frames and a 1000~B minimum in each direction.
Low demand offers one frame per direction; medium offers three. High demand offers
$\max(3,\lfloor L/10\rfloor)$ frames per direction in a contact lasting $L$ seconds.
The service threshold is unchanged. This differs from the older MODERATE constructor that
raises both offered traffic and the minimum; the final derivation changes only offer leaves.
Original nanosecond timestamps and integer tokens are preserved without binary-float
round trips. Every derived model has a leaf-difference audit.

The input-only descriptor report records shortest windows of 58.977355957, 122.687988281
and 74.815063477~s for ISS, NOAA and HST. All exceed the simple 26-second prepared-entry
chain, and no window lies between the corresponding installed/prepared-entry thresholds
of 22 and 26~s. There are one, five and three adjacent A--B--A triples. All 29, 22 and 35
same-station recurrence gaps respectively exceed the base 5400-second cryptoperiod.
For previously activated material, $d_k\leq a+5400\leq e+5400<s$ for prior contact end $e$
and return start $s$. Such material cannot be reused on return. Unactivated staged candidates
remain possible, and this deduction excludes the 10800-second sensitivity. Neither timing
descriptor alone proves service feasibility or general uselessness of an extra horizon.

For this base core, a separate counting consequence is $N_{\rm new}\geq2N_{\rm served}$.
Choose a delivered uplink key and a delivered downlink key for every served contact; both
directional minima are positive. Immutable context/direction bindings prevent sharing across
stations or directions, and the expiry inequality prevents reuse in another same-context
contact. All selected keys are therefore distinct. Because the initial key maps are empty,
each was generated and counted during the execution. Staged unactivated material may be
created earlier but cannot evade that count. Abandoned preparations only add keys. The
source-level argument and exact input premises are recorded in
\texttt{docs/final-research/core-key-lower-bound.md}. This is solely a key-generation bound;
it proves neither useful-time nor CPU/byte optimality and excludes CP10800 sensitivity.

All primary cells use a 5400~s cryptoperiod. FixedPeriod uses $T=3600$~s after successful
SWITCH, and coalesces delayed requests. This is one declared comparator, not an optimally
tuned family of periodic policies. PerContact uses just-in-time preparation. Seed H1 and H2
disable DFS; neither should be confused with the SEED\_ONLY implementation variant that
still searches. The main matched horizon attribution compares Next H1 with Multi H2 using
the same portfolio and limits.

The nine sensitivity attempts change only medium ISS: INSTALL phase durations by 0.5 or 2,
or cryptoperiod to 2700 or 10800~s, for Next/Multi; one additional Multi H3 cell uses the base
model. INSTALL byte counts are unchanged, so this intervention changes modeled CPU/radio
timing rather than uniformly scaling every cost. H3 is explicitly compared with the H2 core
baseline and is not counted as an additional H1/H2 pair.

\section{Decision and Timing Contract}
Inside a contact, H denotes future contacts beyond the current one. Between contacts, it
selects the H nearest future contacts. A search keeps its visible contacts and endpoint fixed
throughout that invocation; entering a later contact inside a hypothetical branch never
reveals a hidden farther contact. Future attacks and outages are unavailable to planning.

The objective first maximizes served contacts and useful secure time, then minimizes
modeled operation CPU, management bytes and new directional keys. Already accrued costs,
delivered progress, counters and issued IDs are preserved. A full-plan rule resolves only
complete equality of all five scores. This does not minimize compromise probability.

All final searches have 5000-node and 1000-ms limits; the full decision has 2000~ms. Seed and
retained-continuation stages each have a 350-ms sublimit and 10000-transition cap. The
external watchdog is 600~s per process. There is no cumulative 20-second controller cutoff.
One heavy JVM runs at a time with a 1024-MiB heap. No budget is increased after a failure.

Offline computation does not advance the physical timeline. Measured mode uses zero fixed
delay plus measured host observation/generation/selection time. This includes seed, replay,
search and selection but excludes separately reported checkpoint I/O and actual execution.
Committed work and actual events continue during delay. A reached or changed WAIT target
causes charged replanning; stale intents are not moved silently. This is one host-dependent
critical-path assumption, not a flight processor benchmark.

Harness controller wall time stops when the controller returns. Its aggregate thread CPU/GC
sample is taken later, after the final checkpoint and most export serialization. These fields
therefore do not define a matched CPU/wall interval. Process wall time includes startup and
exports; individual search telemetry is separate. Checkpoint durations overlap process/controller
time and are not additive overhead on top of either. The detailed boundary table is
\texttt{docs/final-research/timing-definitions.md}; no unavailable export CPU is estimated.

The available-candidate portfolio guarantee is local: all candidates are evaluated from the
same observed root to the same endpoint. The chosen full continuation is no worse under
the unchanged comparator than the evaluated seed, retained suffix or search candidate.
This neither certifies an incomplete search nor ensures that Multi dominates Next across
replanning events. Total searches come from the practical-search ledger; a compact summary
of certified selections omits bounded RULE selections and is not the denominator.

\section{Security, Recovery and Interruption Metrics}
The recording attacker learns one selected active traffic-key pair at an exogenous event.
Replacement material is assumed independent; generators and management credentials are
not compromised. Recorded past and later emitted traffic under the disclosed keys are
retrospectively decryptable. Confidential delivered traffic excludes delivered frames under
keys disclosed by experiment end. Emitted exposure includes aborted prefixes and therefore
cannot be subtracted from delivered bytes as though both totals had the same population.

Four recovery milestones are measured from GS receipt of notice:
\begin{enumerate}
\item confirmed delivery of the revocation notice to the spacecraft;
\item mutually confirmed activation of admissible independent directional keys;
\item actual opening of an admissible SA on independent keys;
\item first delivered independent complete frame.
\end{enumerate}
Independence means exclusion of the disclosed key IDs under the independent-material
assumption, not a statistical test of a real generator. Admissibility checks original material
age, activation deadline, directional binding, revocation and counters. A state already ready
at notice is marked explicitly. Missing events are not replaced by zero. Readiness intervals
also record expiry, binding changes and session closure; first readiness does not imply
continuous readiness. Waiting for key material, notice, contact/channel and offered traffic
can overlap, and their durations are not summed as disjoint causal components.

Channel DOWN aborts committed frames and online primitives with consumed partial costs.
Already admitted aborted frames consume their nonce once. UP does not confirm an interrupted
operation or clear uncertain peer state. INSTALL may be retried when admissible; uncertain
SWITCH/OPEN/REUSE is conservatively blocked without reconciliation. The preregistered ISS
loss is one to three seconds after first AOS; its actual intercepted operation is reported
without shifting the event to obtain an intended result.

\section{Audit, Coverage and Reproduction}
The independent Python auditor checks event ordering, operation phase integrals, complete
and aborted frames, nonce conservation, deadlines, channel/radio overlap, final metrics,
checkpoint/export equality and full-comparator candidate selection. It verifies configuration
and input/JAR identities against registration. Failed processes retain their last usable
checkpoint; partial scores are diagnostic, not final policy quality. A completed physical
execution can contain many incomplete searches. Censoring in a completed incident case,
an unfinished execution and an inapplicable no-target event use distinct labels.

\subsection{Complete Registered Outcome Matrix}

All 108 registered attempts appear below, including process failures and unfinished executions. Case numbers 001--108 are the exact numeric prefixes of the full case IDs in \path{docs/final-research/runs.csv}; that table maps each number to its full ID, frozen model, rule, horizon and invocation identity. Group labels retain block, scenario, load, execution mode and tie rule. Fixed = FixedPeriod; Contact = PerContact. H denotes the registered horizon.

Execution codes: F = FINISHED; PT = PROCESS\_\allowbreak{}TIMEOUT. Codes retain distinct process and controller outcomes.

F means completed physical execution, not a global optimality certificate. I means unavailable final quality for an unfinished execution; NA means the indicated measurement was not exported/available. Neither is zero. Seconds are rounded to two decimal places; positive values below 0.005 s are shown as such. Counts and bytes are integers.

The five objective components are served contacts, useful secure transmission seconds, modeled lifecycle-operation CPU seconds, management bytes, and new directional keys. Controller time is host wall time over the controller call, not the modeled operation CPU.

\begingroup\footnotesize
\setlength{\tabcolsep}{2pt}\renewcommand{\arraystretch}{1.12}
\begin{longtable}{@{}>{\raggedleft\arraybackslash}p{8mm}>{\raggedright\arraybackslash}p{20mm}>{\centering\arraybackslash}p{10mm}>{\raggedleft\arraybackslash}p{10mm}>{\raggedleft\arraybackslash}p{18mm}>{\raggedleft\arraybackslash}p{16mm}>{\raggedleft\arraybackslash}p{20mm}>{\raggedleft\arraybackslash}p{10mm}>{\raggedleft\arraybackslash}p{18mm}@{}}
\caption{All registered outcomes}\label{tab:supp-all-outcomes}\\
\toprule
Case & Policy & Exec. & Served & Useful\newline s & Oper. CPU\newline s & Mgmt.\newline B & Keys & Controller\newline s \\
\midrule\endfirsthead
\multicolumn{9}{l}{\textit{Continued: All registered outcomes}}\\
\toprule
Case & Policy & Exec. & Served & Useful\newline s & Oper. CPU\newline s & Mgmt.\newline B & Keys & Controller\newline s \\
\midrule\endhead
\midrule\multicolumn{9}{r}{\textit{Continued on next page}}\\\endfoot
\bottomrule\endlastfoot
\addlinespace[2pt]\multicolumn{9}{@{}p{149mm}@{}}{\textbf{regression / X5 / historical / offline / CANONICAL}}\\*
001 & Next H1 & F & 2 & 38 & 0.12 & 160 & 0 & 2.34 \\
002 & Multi H2 & F & 2 & 38 & 0.12 & 160 & 0 & 19.63 \\
\addlinespace[2pt]\multicolumn{9}{@{}p{149mm}@{}}{\textbf{regression / X5 / historical / offline / CONTINUE\_\allowbreak{}WAIT}}\\*
003 & Next H1 & F & 2 & 38 & 0.12 & 160 & 0 & 2.36 \\
004 & Multi H2 & F & 2 & 38 & 0.12 & 160 & 0 & 21.2 \\
\addlinespace[2pt]\multicolumn{9}{@{}p{149mm}@{}}{\textbf{regression / Y / historical / offline / CANONICAL}}\\*
005 & Next H1 & F & 3 & 42 & 0.18 & 240 & 0 & 2.78 \\
006 & Multi H2 & F & 3 & 42 & 0.18 & 240 & 0 & 2.98 \\
\addlinespace[2pt]\multicolumn{9}{@{}p{149mm}@{}}{\textbf{regression / Y / historical / offline / CONTINUE\_\allowbreak{}WAIT}}\\*
007 & Next H1 & F & 3 & 42 & 0.18 & 240 & 0 & 2.84 \\
008 & Multi H2 & F & 3 & 42 & 0.18 & 240 & 0 & 2.8 \\
\addlinespace[2pt]\multicolumn{9}{@{}p{149mm}@{}}{\textbf{regression / demand-none / historical / offline / CANONICAL}}\\*
009 & Next H1 & F & 3 & 410 & 2.13 & 3240 & 6 & 2.94 \\
010 & Multi H2 & F & 3 & 410 & 2.13 & 3240 & 6 & 17.9 \\
\addlinespace[2pt]\multicolumn{9}{@{}p{149mm}@{}}{\textbf{regression / demand-none / historical / offline / CONTINUE\_\allowbreak{}WAIT}}\\*
011 & Next H1 & F & 3 & 410 & 2.13 & 3240 & 6 & 2.14 \\
012 & Multi H2 & F & 3 & 410 & 2.13 & 3240 & 6 & 5.75 \\
\addlinespace[2pt]\multicolumn{9}{@{}p{149mm}@{}}{\textbf{regression / X5 / historical / measured / CANONICAL}}\\*
013 & Next H1 & F & 2 & 33.75 & 0.12 & 160 & 0 & 5.61 \\
014 & Multi H2 & F & 2 & 27.5 & 0.12 & 160 & 0 & 12.12 \\
\addlinespace[2pt]\multicolumn{9}{@{}p{149mm}@{}}{\textbf{regression / Y / historical / measured / CANONICAL}}\\*
015 & Next H1 & F & 3 & 37 & 0.18 & 240 & 0 & 6.4 \\
016 & Multi H2 & F & 3 & 37 & 0.18 & 240 & 0 & 6.63 \\
\addlinespace[2pt]\multicolumn{9}{@{}p{149mm}@{}}{\textbf{regression / demand-none / historical / measured / CANONICAL}}\\*
017 & Next H1 & F & 3 & 410 & 2.13 & 3240 & 6 & 3.98 \\
018 & Multi H2 & F & 3 & 410 & 2.13 & 3240 & 6 & 7.33 \\
\addlinespace[2pt]\multicolumn{9}{@{}p{149mm}@{}}{\textbf{core / ISS46 / low / offline / CANONICAL}}\\*
019 & Fixed & F & 46 & 920 & 35.66 & 49680 & 122 & 7.56 \\
020 & Contact & F & 46 & 920 & 32.66 & 49680 & 92 & 6.59 \\
021 & Next H1 & F & 46 & 920 & 32.66 & 49680 & 92 & 213.28 \\
022 & Multi H2 & F & 46 & 920 & 32.66 & 49680 & 92 & 521.23 \\
023 & Seed H1 & F & 46 & 920 & 32.66 & 49680 & 92 & 16.29 \\
024 & Seed H2 & F & 46 & 920 & 32.66 & 49680 & 92 & 23.17 \\
\addlinespace[2pt]\multicolumn{9}{@{}p{149mm}@{}}{\textbf{core / ISS46 / medium / offline / CANONICAL}}\\*
025 & Fixed & F & 46 & 2740 & 35.66 & 49680 & 122 & 11.72 \\
026 & Contact & F & 46 & 2740 & 32.66 & 49680 & 92 & 10.04 \\
027 & Next H1 & F & 46 & 2740 & 32.66 & 49680 & 92 & 527.78 \\
028 & Multi H2 & PT & I & I & I & I & I & NA \\
029 & Seed H1 & F & 46 & 2740 & 32.66 & 49680 & 92 & 25.38 \\
030 & Seed H2 & F & 46 & 2740 & 32.66 & 49680 & 92 & 35.94 \\
\addlinespace[2pt]\multicolumn{9}{@{}p{149mm}@{}}{\textbf{core / ISS46 / high / offline / CANONICAL}}\\*
031 & Fixed & F & 46 & 14390 & 35.66 & 49680 & 122 & 37.7 \\
032 & Contact & F & 46 & 14390 & 32.66 & 49680 & 92 & 33.68 \\
033 & Next H1 & PT & I & I & I & I & I & NA \\
034 & Multi H2 & PT & I & I & I & I & I & NA \\
035 & Seed H1 & F & 46 & 14390 & 32.66 & 49680 & 92 & 132.68 \\
036 & Seed H2 & F & 46 & 14390 & 32.66 & 49680 & 92 & 191.42 \\
\addlinespace[2pt]\multicolumn{9}{@{}p{149mm}@{}}{\textbf{core / NOAA38 / low / offline / CANONICAL}}\\*
037 & Fixed & F & 38 & 760 & 29.78 & 41040 & 104 & 6.82 \\
038 & Contact & F & 38 & 760 & 26.98 & 41040 & 76 & 5.59 \\
039 & Next H1 & F & 38 & 760 & 26.98 & 41040 & 76 & 139.61 \\
040 & Multi H2 & F & 38 & 760 & 26.98 & 41040 & 76 & 389.14 \\
041 & Seed H1 & F & 38 & 760 & 26.98 & 41040 & 76 & 11.53 \\
042 & Seed H2 & F & 38 & 760 & 26.98 & 41040 & 76 & 15.84 \\
\addlinespace[2pt]\multicolumn{9}{@{}p{149mm}@{}}{\textbf{core / NOAA38 / medium / offline / CANONICAL}}\\*
043 & Fixed & F & 38 & 2280 & 29.78 & 41040 & 104 & 8.61 \\
044 & Contact & F & 38 & 2280 & 26.98 & 41040 & 76 & 7.24 \\
045 & Next H1 & F & 38 & 2280 & 26.98 & 41040 & 76 & 337.17 \\
046 & Multi H2 & F & 38 & 2280 & 26.98 & 41040 & 76 & 569.84 \\
047 & Seed H1 & F & 38 & 2280 & 26.98 & 41040 & 76 & 17.86 \\
048 & Seed H2 & F & 38 & 2280 & 26.98 & 41040 & 76 & 24.92 \\
\addlinespace[2pt]\multicolumn{9}{@{}p{149mm}@{}}{\textbf{core / NOAA38 / high / offline / CANONICAL}}\\*
049 & Fixed & F & 38 & 19210 & 29.78 & 41040 & 104 & 42 \\
050 & Contact & F & 38 & 19210 & 26.98 & 41040 & 76 & 40.69 \\
051 & Next H1 & PT & I & I & I & I & I & NA \\
052 & Multi H2 & PT & I & I & I & I & I & NA \\
053 & Seed H1 & F & 38 & 19210 & 26.98 & 41040 & 76 & 168.14 \\
054 & Seed H2 & F & 38 & 19210 & 26.98 & 41040 & 76 & 251.68 \\
\addlinespace[2pt]\multicolumn{9}{@{}p{149mm}@{}}{\textbf{core / HST48 / low / offline / CANONICAL}}\\*
055 & Fixed & F & 48 & 960 & 36.28 & 51840 & 118 & 8.96 \\
056 & Contact & F & 48 & 960 & 34.08 & 51840 & 96 & 7.52 \\
057 & Next H1 & F & 48 & 960 & 34.08 & 51840 & 96 & 97.57 \\
058 & Multi H2 & F & 48 & 960 & 34.08 & 51840 & 96 & 453.39 \\
059 & Seed H1 & F & 48 & 960 & 34.08 & 51840 & 96 & 15.46 \\
060 & Seed H2 & F & 48 & 960 & 34.08 & 51840 & 96 & 21.58 \\
\addlinespace[2pt]\multicolumn{9}{@{}p{149mm}@{}}{\textbf{core / HST48 / medium / offline / CANONICAL}}\\*
061 & Fixed & F & 48 & 2880 & 36.28 & 51840 & 118 & 12.05 \\
062 & Contact & F & 48 & 2880 & 34.08 & 51840 & 96 & 10.81 \\
063 & Next H1 & F & 48 & 2880 & 34.08 & 51840 & 96 & 319.17 \\
064 & Multi H2 & PT & I & I & I & I & I & NA \\
065 & Seed H1 & F & 48 & 2880 & 34.08 & 51840 & 96 & 23.63 \\
066 & Seed H2 & F & 48 & 2880 & 34.08 & 51840 & 96 & 34.31 \\
\addlinespace[2pt]\multicolumn{9}{@{}p{149mm}@{}}{\textbf{core / HST48 / high / offline / CANONICAL}}\\*
067 & Fixed & F & 48 & 15560 & 36.28 & 51840 & 118 & 37.23 \\
068 & Contact & F & 48 & 15560 & 34.08 & 51840 & 96 & 35.35 \\
069 & Next H1 & PT & I & I & I & I & I & NA \\
070 & Multi H2 & PT & I & I & I & I & I & NA \\
071 & Seed H1 & F & 48 & 15560 & 34.08 & 51840 & 96 & 118.73 \\
072 & Seed H2 & F & 48 & 15560 & 34.08 & 51840 & 96 & 172.56 \\
\addlinespace[2pt]\multicolumn{9}{@{}p{149mm}@{}}{\textbf{measured / ISS46 / low / measured / CANONICAL}}\\*
073 & Fixed & F & 46 & 920 & 35.66 & 49680 & 122 & 9.38 \\
074 & Contact & F & 46 & 920 & 32.66 & 49680 & 92 & 8.55 \\
075 & Next H1 & F & 46 & 920 & 32.66 & 49680 & 92 & 245.48 \\
076 & Multi H2 & F & 46 & 920 & 32.66 & 49680 & 92 & 544 \\
077 & Seed H1 & F & 46 & 920 & 32.66 & 49680 & 92 & 17.91 \\
\addlinespace[2pt]\multicolumn{9}{@{}p{149mm}@{}}{\textbf{measured / ISS46 / medium / measured / CANONICAL}}\\*
078 & Fixed & F & 46 & 2740 & 35.66 & 49680 & 122 & 12.44 \\
079 & Contact & F & 46 & 2740 & 32.66 & 49680 & 92 & 11.22 \\
080 & Next H1 & F & 46 & 2720 & 32.66 & 49680 & 92 & 517.06 \\
081 & Multi H2 & PT & I & I & I & I & I & NA \\
082 & Seed H1 & F & 46 & 2740 & 32.66 & 49680 & 92 & 28.35 \\
\addlinespace[2pt]\multicolumn{9}{@{}p{149mm}@{}}{\textbf{measured / ISS46 / high / measured / CANONICAL}}\\*
083 & Fixed & F & 46 & 14390 & 35.66 & 49680 & 122 & 36.29 \\
084 & Contact & F & 46 & 14390 & 32.66 & 49680 & 92 & 34.64 \\
085 & Next H1 & PT & I & I & I & I & I & NA \\
086 & Multi H2 & PT & I & I & I & I & I & NA \\
087 & Seed H1 & F & 46 & 14290 & 32.66 & 49680 & 92 & 132.28 \\
\addlinespace[2pt]\multicolumn{9}{@{}p{149mm}@{}}{\textbf{sensitivity / ISS46 / medium / offline / CANONICAL; factor/model ISS46-medium-install0.5}}\\*
088 & Next H1 & F & 46 & 2740 & 25.76 & 49680 & 92 & 514.33 \\
089 & Multi H2 & PT & I & I & I & I & I & NA \\
\addlinespace[2pt]\multicolumn{9}{@{}p{149mm}@{}}{\textbf{sensitivity / ISS46 / medium / offline / CANONICAL; factor/model ISS46-medium-install2}}\\*
090 & Next H1 & F & 46 & 2730 & 46.46 & 49680 & 92 & 509.22 \\
091 & Multi H2 & PT & I & I & I & I & I & NA \\
\addlinespace[2pt]\multicolumn{9}{@{}p{149mm}@{}}{\textbf{sensitivity / ISS46 / medium / offline / CANONICAL; factor/model ISS46-medium-cp2700}}\\*
092 & Next H1 & F & 46 & 2740 & 32.66 & 49680 & 92 & 506.22 \\
093 & Multi H2 & PT & I & I & I & I & I & NA \\
\addlinespace[2pt]\multicolumn{9}{@{}p{149mm}@{}}{\textbf{sensitivity / ISS46 / medium / offline / CANONICAL; factor/model ISS46-medium-cp10800}}\\*
094 & Next H1 & F & 46 & 2740 & 22.91 & 34680 & 62 & 266.51 \\
095 & Multi H2 & F & 46 & 2740 & 22.91 & 34680 & 62 & 462.1 \\
\addlinespace[2pt]\multicolumn{9}{@{}p{149mm}@{}}{\textbf{sensitivity / ISS46 / medium / offline / CANONICAL; factor/model ISS46-medium}}\\*
096 & Multi H3 & PT & I & I & I & I & I & NA \\
\addlinespace[2pt]\multicolumn{9}{@{}p{149mm}@{}}{\textbf{security / X-undetected / historical / offline / CANONICAL}}\\*
097 & Next H1 & F & 2 & 38 & 0.12 & 160 & 0 & 3.07 \\
098 & Multi H2 & F & 2 & 38 & 0.12 & 160 & 0 & 24.03 \\
\addlinespace[2pt]\multicolumn{9}{@{}p{149mm}@{}}{\textbf{security / X-detected2 / historical / offline / CANONICAL}}\\*
099 & Next H1 & F & 3 & 33.5 & 0.93 & 1416 & 2 & 3.35 \\
100 & Multi H2 & F & 3 & 33.5 & 0.93 & 1416 & 2 & 16.31 \\
\addlinespace[2pt]\multicolumn{9}{@{}p{149mm}@{}}{\textbf{security / X-gap / historical / offline / CANONICAL}}\\*
101 & Next H1 & F & 2 & 38 & 0.66 & 1056 & 2 & 2.73 \\
102 & Multi H2 & F & 2 & 38 & 0.66 & 1056 & 2 & 22.45 \\
\addlinespace[2pt]\multicolumn{9}{@{}p{149mm}@{}}{\textbf{security / X-channel-temporary / historical / offline / CANONICAL}}\\*
103 & Next H1 & F & 3 & 32 & 0.89 & 1320 & 2 & 3.62 \\
104 & Multi H2 & F & 3 & 32 & 0.89 & 1320 & 2 & 10.34 \\
\addlinespace[2pt]\multicolumn{9}{@{}p{149mm}@{}}{\textbf{security / demand-late-detected / historical / offline / CANONICAL}}\\*
105 & Next H1 & F & 3 & 400 & 2.88 & 4416 & 8 & 3.35 \\
106 & Multi H2 & F & 3 & 410 & 2.88 & 4416 & 8 & 18.71 \\
\addlinespace[2pt]\multicolumn{9}{@{}p{149mm}@{}}{\textbf{management / ISS3 / low / offline / CANONICAL}}\\*
107 & Next H1 & F & 3 & 60 & 2.23 & 3600 & 6 & 3.9 \\
108 & Multi H2 & F & 3 & 60 & 2.23 & 3600 & 6 & 9.14 \\
\end{longtable}\endgroup


\subsection{Actual Computational Work}

Calls, incomplete calls, nodes and search seconds come from the exported actual-search ledger, including complete and incomplete invocations. They are not the count of certified selections. A genuine empty ledger can report zero; an unexported ledger after process termination remains NA. Process wall time includes startup and export/termination. Search CPU covers the calling thread during search; it is separate from aggregate harness CPU/GC with its different sampling interval. Miss denotes decision deadline misses. These are single host observations, not onboard timing estimates.

\begingroup\footnotesize
\setlength{\tabcolsep}{2pt}\renewcommand{\arraystretch}{1.12}
\begin{longtable}{@{}>{\raggedleft\arraybackslash}p{8mm}>{\centering\arraybackslash}p{10mm}>{\raggedleft\arraybackslash}p{16mm}>{\raggedleft\arraybackslash}p{12mm}>{\raggedleft\arraybackslash}p{12mm}>{\raggedleft\arraybackslash}p{18mm}>{\raggedleft\arraybackslash}p{18mm}>{\raggedleft\arraybackslash}p{18mm}>{\raggedleft\arraybackslash}p{12mm}@{}}
\caption{Process and search-ledger diagnostics}\label{tab:supp-search}\\
\toprule
Case & Exec. & Process\newline s & Calls & Incomp.\newline calls & Nodes & Search\newline s & Search CPU\newline s & Miss \\
\midrule\endfirsthead
\multicolumn{9}{l}{\textit{Continued: Process and search-ledger diagnostics}}\\
\toprule
Case & Exec. & Process\newline s & Calls & Incomp.\newline calls & Nodes & Search\newline s & Search CPU\newline s & Miss \\
\midrule\endhead
\midrule\multicolumn{9}{r}{\textit{Continued on next page}}\\\endfoot
\bottomrule\endlastfoot
001 & F & 3.03 & 13 & 0 & 3738 & 0.66 & 0.67 & 0 \\
002 & F & 20.41 & 78 & 67 & 340824 & 17.69 & 17.48 & 0 \\
003 & F & 3.09 & 13 & 0 & 3738 & 0.68 & 0.66 & 0 \\
004 & F & 21.97 & 78 & 67 & 340824 & 19.18 & 18.77 & 0 \\
005 & F & 3.55 & 15 & 0 & 4215 & 0.82 & 0.75 & 0 \\
006 & F & 3.73 & 15 & 0 & 4932 & 0.93 & 0.88 & 0 \\
007 & F & 3.66 & 15 & 0 & 4215 & 0.82 & 0.81 & 0 \\
008 & F & 3.56 & 15 & 0 & 4932 & 0.91 & 0.91 & 0 \\
009 & F & 3.8 & 32 & 0 & 7878 & 1.67 & 1.59 & 0 \\
010 & F & 18.67 & 68 & 48 & 241778 & 16.8 & 16.41 & 0 \\
011 & F & 2.94 & 33 & 0 & 3390 & 0.86 & 0.88 & 0 \\
012 & F & 6.48 & 33 & 8 & 41546 & 4.86 & 4.8 & 0 \\
013 & F & 6.56 & 148 & 0 & 43379 & 3.3 & 3.27 & 0 \\
014 & F & 13.08 & 124 & 31 & 177633 & 10.59 & 10.48 & 0 \\
015 & F & 7.45 & 163 & 0 & 47574 & 3.72 & 3.59 & 0 \\
016 & F & 7.69 & 163 & 0 & 52701 & 4.04 & 3.86 & 0 \\
017 & F & 5.16 & 74 & 1 & 15845 & 2.67 & 2.56 & 0 \\
018 & F & 8.22 & 73 & 11 & 61214 & 6.34 & 6.17 & 0 \\
019 & F & 8.76 & 0 & 0 & 0 & 0 & 0 & 0 \\
020 & F & 7.74 & 0 & 0 & 0 & 0 & 0 & 0 \\
021 & F & 214.84 & 570 & 209 & 1019808 & 191.59 & 189.55 & 0 \\
022 & F & 523.12 & 614 & 593 & 2470275 & 484.57 & 478.8 & 0 \\
023 & F & 17.67 & 0 & 0 & 0 & 0 & 0 & 0 \\
024 & F & 24.55 & 0 & 0 & 0 & 0 & 0 & 0 \\
025 & F & 13.36 & 0 & 0 & 0 & 0 & 0 & 0 \\
026 & F & 11.36 & 0 & 0 & 0 & 0 & 0 & 0 \\
027 & F & 529.45 & 734 & 561 & 2278035 & 487.78 & 481.58 & 0 \\
028 & PT & 600.05 & NA & NA & NA & NA & NA & NA \\
029 & F & 26.84 & 0 & 0 & 0 & 0 & 0 & 0 \\
030 & F & 37.61 & 0 & 0 & 0 & 0 & 0 & 0 \\
031 & F & 40.55 & 0 & 0 & 0 & 0 & 0 & 0 \\
032 & F & 35.97 & 0 & 0 & 0 & 0 & 0 & 0 \\
033 & PT & 600.03 & NA & NA & NA & NA & NA & NA \\
034 & PT & 600.05 & NA & NA & NA & NA & NA & NA \\
035 & F & 135.67 & 0 & 0 & 0 & 0 & 0 & 0 \\
036 & F & 195.05 & 0 & 0 & 0 & 0 & 0 & 0 \\
037 & F & 7.94 & 0 & 0 & 0 & 0 & 0 & 0 \\
038 & F & 6.64 & 0 & 0 & 0 & 0 & 0 & 0 \\
039 & F & 140.88 & 456 & 125 & 634298 & 122.54 & 121.47 & 0 \\
040 & F & 390.72 & 513 & 467 & 2139796 & 364.19 & 359.31 & 0 \\
041 & F & 12.72 & 0 & 0 & 0 & 0 & 0 & 0 \\
042 & F & 16.95 & 0 & 0 & 0 & 0 & 0 & 0 \\
043 & F & 10.17 & 0 & 0 & 0 & 0 & 0 & 0 \\
044 & F & 8.31 & 0 & 0 & 0 & 0 & 0 & 0 \\
045 & F & 338.62 & 590 & 391 & 1847650 & 312.15 & 308.58 & 0 \\
046 & F & 571.56 & 663 & 655 & 2874947 & 529.26 & 522.62 & 0 \\
047 & F & 19.09 & 0 & 0 & 0 & 0 & 0 & 0 \\
048 & F & 26.26 & 0 & 0 & 0 & 0 & 0 & 0 \\
049 & F & 44.86 & 0 & 0 & 0 & 0 & 0 & 0 \\
050 & F & 43.05 & 0 & 0 & 0 & 0 & 0 & 0 \\
051 & PT & 600.03 & NA & NA & NA & NA & NA & NA \\
052 & PT & 600.03 & NA & NA & NA & NA & NA & NA \\
053 & F & 172.11 & 0 & 0 & 0 & 0 & 0 & 0 \\
054 & F & 256.33 & 0 & 0 & 0 & 0 & 0 & 0 \\
055 & F & 10.17 & 0 & 0 & 0 & 0 & 0 & 0 \\
056 & F & 8.69 & 0 & 0 & 0 & 0 & 0 & 0 \\
057 & F & 99.03 & 576 & 35 & 386357 & 76.83 & 75.88 & 0 \\
058 & F & 454.97 & 650 & 511 & 2301201 & 421.14 & 416.3 & 0 \\
059 & F & 16.7 & 0 & 0 & 0 & 0 & 0 & 0 \\
060 & F & 22.88 & 0 & 0 & 0 & 0 & 0 & 0 \\
061 & F & 13.59 & 0 & 0 & 0 & 0 & 0 & 0 \\
062 & F & 12.06 & 0 & 0 & 0 & 0 & 0 & 0 \\
063 & F & 320.83 & 655 & 338 & 1581737 & 288.17 & 285.11 & 0 \\
064 & PT & 600.03 & NA & NA & NA & NA & NA & NA \\
065 & F & 25.08 & 0 & 0 & 0 & 0 & 0 & 0 \\
066 & F & 35.86 & 0 & 0 & 0 & 0 & 0 & 0 \\
067 & F & 39.86 & 0 & 0 & 0 & 0 & 0 & 0 \\
068 & F & 37.5 & 0 & 0 & 0 & 0 & 0 & 0 \\
069 & PT & 600.03 & NA & NA & NA & NA & NA & NA \\
070 & PT & 600.03 & NA & NA & NA & NA & NA & NA \\
071 & F & 121.66 & 0 & 0 & 0 & 0 & 0 & 0 \\
072 & F & 176.17 & 0 & 0 & 0 & 0 & 0 & 0 \\
073 & F & 11.31 & 0 & 0 & 0 & 0 & 0 & 0 \\
074 & F & 10.2 & 0 & 0 & 0 & 0 & 0 & 0 \\
075 & F & 247.28 & 636 & 229 & 1089372 & 224.59 & 222.41 & 0 \\
076 & F & 545.92 & 632 & 611 & 2524221 & 516.18 & 510.34 & 0 \\
077 & F & 19.61 & 0 & 0 & 0 & 0 & 0 & 0 \\
078 & F & 14.73 & 0 & 0 & 0 & 0 & 0 & 0 \\
079 & F & 13.05 & 0 & 0 & 0 & 0 & 0 & 0 \\
080 & F & 519.24 & 819 & 573 & 2481749 & 486.49 & 480.64 & 0 \\
081 & PT & 600.05 & NA & NA & NA & NA & NA & NA \\
082 & F & 30.36 & 0 & 0 & 0 & 0 & 0 & 0 \\
083 & F & 40.42 & 0 & 0 & 0 & 0 & 0 & 0 \\
084 & F & 37.98 & 0 & 0 & 0 & 0 & 0 & 0 \\
085 & PT & 600.03 & NA & NA & NA & NA & NA & NA \\
086 & PT & 600.03 & NA & NA & NA & NA & NA & NA \\
087 & F & 136.97 & 0 & 0 & 0 & 0 & 0 & 0 \\
088 & F & 516.03 & 733 & 555 & 2330193 & 476.45 & 470.44 & 0 \\
089 & PT & 600.05 & NA & NA & NA & NA & NA & NA \\
090 & F & 510.92 & 729 & 551 & 2310270 & 472.21 & 465.64 & 0 \\
091 & PT & 600.05 & NA & NA & NA & NA & NA & NA \\
092 & F & 507.86 & 734 & 553 & 2410550 & 469.35 & 463.28 & 0 \\
093 & PT & 600.05 & NA & NA & NA & NA & NA & NA \\
094 & F & 267.92 & 551 & 312 & 1571855 & 245.24 & 242.47 & 0 \\
095 & F & 463.72 & 636 & 537 & 2493989 & 427.87 & 422.55 & 0 \\
096 & PT & 600.03 & NA & NA & NA & NA & NA & NA \\
097 & F & 3.92 & 13 & 0 & 3738 & 0.94 & 0.86 & 0 \\
098 & F & 25.06 & 78 & 67 & 340680 & 21.47 & 21.23 & 0 \\
099 & F & 4.16 & 20 & 0 & 6210 & 1.16 & 1.09 & 0 \\
100 & F & 17.55 & 55 & 38 & 195794 & 13.78 & 13.59 & 0 \\
101 & F & 3.53 & 15 & 0 & 3764 & 0.78 & 0.73 & 0 \\
102 & F & 23.33 & 80 & 67 & 340293 & 20.02 & 19.88 & 0 \\
103 & F & 4.44 & 24 & 0 & 8405 & 1.43 & 1.34 & 0 \\
104 & F & 11.14 & 39 & 18 & 99999 & 8.23 & 8.12 & 0 \\
105 & F & 4.14 & 38 & 1 & 8054 & 1.85 & 1.89 & 0 \\
106 & F & 19.5 & 73 & 45 & 230766 & 17.36 & 16.98 & 0 \\
107 & F & 4.69 & 35 & 0 & 18224 & 2.73 & 2.59 & 0 \\
108 & F & 9.89 & 35 & 14 & 77606 & 8.1 & 7.98 & 0 \\
\end{longtable}\endgroup


Measured decision host seconds sum the actual decision intervals charged to the model. They exclude checkpoint I/O and actual executor work and are not interchangeable with controller wall time. Recorded diagnostics can describe a stopped execution; they do not supply its unavailable final service result.

\begingroup\footnotesize
\setlength{\tabcolsep}{2pt}\renewcommand{\arraystretch}{1.12}
\begin{longtable}{@{}>{\raggedleft\arraybackslash}p{12mm}>{\centering\arraybackslash}p{14mm}>{\raggedleft\arraybackslash}p{22mm}>{\raggedleft\arraybackslash}p{35mm}>{\raggedleft\arraybackslash}p{22mm}@{}}
\caption{Measured decision-delay diagnostics}\label{tab:supp-delay}\\
\toprule
Case & Exec. & Decisions & Decision host\newline s & Deadline\newline misses \\
\midrule\endfirsthead
\multicolumn{5}{l}{\textit{Continued: Measured decision-delay diagnostics}}\\
\toprule
Case & Exec. & Decisions & Decision host\newline s & Deadline\newline misses \\
\midrule\endhead
\midrule\multicolumn{5}{r}{\textit{Continued on next page}}\\\endfoot
\bottomrule\endlastfoot
013 & F & 148 & 4.62 & 0 \\
014 & F & 124 & 11.39 & 0 \\
015 & F & 163 & 5.22 & 0 \\
016 & F & 163 & 5.54 & 0 \\
017 & F & 74 & 3.24 & 0 \\
018 & F & 73 & 6.77 & 0 \\
073 & F & 713 & 0.78 & 0 \\
074 & F & 639 & 0.94 & 0 \\
075 & F & 636 & 238.54 & 0 \\
076 & F & 632 & 537.02 & 0 \\
077 & F & 639 & 10.74 & 0 \\
078 & F & 894 & 0.99 & 0 \\
079 & F & 821 & 1.15 & 0 \\
080 & F & 819 & 507.56 & 0 \\
081 & PT & NA & NA & NA \\
082 & F & 821 & 17.65 & 0 \\
083 & F & 2060 & 2.08 & 0 \\
084 & F & 1986 & 2.74 & 0 \\
085 & PT & NA & NA & NA \\
086 & PT & NA & NA & NA \\
087 & F & 1975 & 101.23 & 0 \\
\end{longtable}\endgroup


\subsection{Disclosure and Interruption Outcomes}

The following tables retain every registered security and management-interruption case. UL/DL denote uplink/downlink. Delivered bytes and frame counts include complete delivered frames only. Confidential bytes are retrospective through experiment end. Decryptable emitted bytes include recorded traffic before disclosure and traffic emitted afterward, including aborted prefixes. Exposure and confidential delivery are therefore not a partition of delivered traffic. Only FINISHED executions receive final traffic totals.

\begingroup\footnotesize
\setlength{\tabcolsep}{2pt}\renewcommand{\arraystretch}{1.12}
\begin{longtable}{@{}>{\raggedleft\arraybackslash}p{8mm}>{\centering\arraybackslash}p{10mm}>{\raggedleft\arraybackslash}p{30mm}>{\raggedleft\arraybackslash}p{24mm}>{\raggedleft\arraybackslash}p{30mm}>{\raggedleft\arraybackslash}p{30mm}@{}}
\caption{Final delivered traffic and emitted exposure}\label{tab:supp-traffic}\\
\toprule
Case & Exec. & Delivered B\newline UL/DL & Frames\newline UL/DL & Confidential B\newline UL/DL & Decryptable B\newline before/after \\
\midrule\endfirsthead
\multicolumn{6}{l}{\textit{Continued: Final delivered traffic and emitted exposure}}\\
\toprule
Case & Exec. & Delivered B\newline UL/DL & Frames\newline UL/DL & Confidential B\newline UL/DL & Decryptable B\newline before/after \\
\midrule\endhead
\midrule\multicolumn{6}{r}{\textit{Continued on next page}}\\\endfoot
\bottomrule\endlastfoot
097 & F & 7600/\allowbreak{}7600 & 76/\allowbreak{}76 & 3800/\allowbreak{}3800 & 2800/\allowbreak{}4800 \\
098 & F & 7600/\allowbreak{}7600 & 76/\allowbreak{}76 & 3800/\allowbreak{}3800 & 2800/\allowbreak{}4800 \\
099 & F & 6700/\allowbreak{}6700 & 67/\allowbreak{}67 & 4900/\allowbreak{}4900 & 2800/\allowbreak{}800 \\
100 & F & 6700/\allowbreak{}6700 & 67/\allowbreak{}67 & 4900/\allowbreak{}4900 & 2800/\allowbreak{}800 \\
101 & F & 7600/\allowbreak{}7600 & 76/\allowbreak{}76 & 3800/\allowbreak{}3800 & 7600/\allowbreak{}0 \\
102 & F & 7600/\allowbreak{}7600 & 76/\allowbreak{}76 & 3800/\allowbreak{}3800 & 7600/\allowbreak{}0 \\
103 & F & 6400/\allowbreak{}6400 & 64/\allowbreak{}64 & 6400/\allowbreak{}6400 & 0/\allowbreak{}0 \\
104 & F & 6400/\allowbreak{}6400 & 64/\allowbreak{}64 & 6400/\allowbreak{}6400 & 0/\allowbreak{}0 \\
105 & F & 20000/\allowbreak{}20000 & 20/\allowbreak{}20 & 3000/\allowbreak{}3000 & 33842/\allowbreak{}200 \\
106 & F & 21000/\allowbreak{}20000 & 21/\allowbreak{}20 & 4000/\allowbreak{}3000 & 33842/\allowbreak{}158 \\
107 & F & 3000/\allowbreak{}3000 & 3/\allowbreak{}3 & 3000/\allowbreak{}3000 & 0/\allowbreak{}0 \\
108 & F & 3000/\allowbreak{}3000 & 3/\allowbreak{}3 & 3000/\allowbreak{}3000 & 0/\allowbreak{}0 \\
\end{longtable}\endgroup


Milestones are seconds after GS notice: E = observed event, AR = already ready at the origin, C = right-censored in a completed execution, U = unavailable origin or missing evidence, N/A = not applicable (no attack, no target or undetected incident). I retains an unfinished execution without presenting its prefix as final recovery or final censoring. Incident state is reported separately. First readiness does not imply continuous availability.

\begingroup\footnotesize
\setlength{\tabcolsep}{2pt}\renewcommand{\arraystretch}{1.12}
\begin{longtable}{@{}>{\raggedleft\arraybackslash}p{8mm}>{\centering\arraybackslash}p{10mm}>{\raggedright\arraybackslash}p{30mm}>{\raggedleft\arraybackslash}p{22mm}>{\raggedleft\arraybackslash}p{22mm}>{\raggedleft\arraybackslash}p{22mm}>{\raggedleft\arraybackslash}p{22mm}@{}}
\caption{Four distinct recovery milestones}\label{tab:supp-recovery}\\
\toprule
Case & Exec. & Incident state & SC notice\newline s & Indep. keys\newline s & Open SA\newline s & Indep. frame\newline s \\
\midrule\endfirsthead
\multicolumn{7}{l}{\textit{Continued: Four distinct recovery milestones}}\\
\toprule
Case & Exec. & Incident state & SC notice\newline s & Indep. keys\newline s & Open SA\newline s & Indep. frame\newline s \\
\midrule\endhead
\midrule\multicolumn{7}{r}{\textit{Continued on next page}}\\\endfoot
\bottomrule\endlastfoot
097 & F & UNDETECTED & N/A & N/A & N/A & N/A \\
098 & F & UNDETECTED & N/A & N/A & N/A & N/A \\
099 & F & RECOVERED & E:0.4 & E:7.4 & E:8.4 & E:8.65 \\
100 & F & RECOVERED & E:0.4 & E:7.4 & E:8.4 & E:8.65 \\
101 & F & CENSORED\_\allowbreak{}NO\_\allowbreak{}RECOVERY & E:18.4 & C & C & C \\
102 & F & CENSORED\_\allowbreak{}NO\_\allowbreak{}RECOVERY & E:18.4 & C & C & C \\
103 & F & NO\_\allowbreak{}ATTACK & N/A & N/A & N/A & N/A \\
104 & F & NO\_\allowbreak{}ATTACK & N/A & N/A & N/A & N/A \\
105 & F & RECOVERED & E:0.4 & E:7.4 & E:8.4 & E:18.4 \\
106 & F & RECOVERED & E:3.98 & E:4.98 & E:5.98 & E:15.98 \\
107 & F & NO\_\allowbreak{}ATTACK & N/A & N/A & N/A & N/A \\
108 & F & NO\_\allowbreak{}ATTACK & N/A & N/A & N/A & N/A \\
\end{longtable}\endgroup


Interruption diagnostics below describe actually recorded events and consumed partial operation resources, not unused budgets or additional costs to add to the final five scores. For an unfinished execution they remain prefix diagnostics; absent exports remain NA. The fixed outage is not shifted to target a different operation. Full action counts remain in runs.csv.

\begingroup\footnotesize
\setlength{\tabcolsep}{2pt}\renewcommand{\arraystretch}{1.12}
\begin{longtable}{@{}>{\raggedleft\arraybackslash}p{8mm}>{\centering\arraybackslash}p{10mm}>{\raggedright\arraybackslash}p{53mm}>{\raggedleft\arraybackslash}p{18mm}>{\raggedleft\arraybackslash}p{23mm}>{\raggedleft\arraybackslash}p{23mm}@{}}
\caption{Observed interruption diagnostics}\label{tab:supp-interruption}\\
\toprule
Case & Exec. & Interrupted actions/counts & Aborted\newline frames & Partial oper.\newline CPU s & Partial mgmt.\newline B \\
\midrule\endfirsthead
\multicolumn{6}{l}{\textit{Continued: Observed interruption diagnostics}}\\
\toprule
Case & Exec. & Interrupted actions/counts & Aborted\newline frames & Partial oper.\newline CPU s & Partial mgmt.\newline B \\
\midrule\endhead
\midrule\multicolumn{6}{r}{\textit{Continued on next page}}\\\endfoot
\bottomrule\endlastfoot
097 & F & \{\} & 0 & 0 & 0 \\
098 & F & \{\} & 0 & 0 & 0 \\
099 & F & \{\} & 0 & 0 & 0 \\
100 & F & \{\} & 0 & 0 & 0 \\
101 & F & \{\} & 0 & 0 & 0 \\
102 & F & \{\} & 0 & 0 & 0 \\
103 & F & \{\} & 0 & 0 & 0 \\
104 & F & \{\} & 0 & 0 & 0 \\
105 & F & \{\} & 1 & 0 & 0 \\
106 & F & \{\} & 0 & 0 & 0 \\
107 & F & \{"INSTALL": 1\} & 0 & 0.1 & 360 \\
108 & F & \{"INSTALL": 1\} & 0 & 0.1 & 360 \\
\end{longtable}\endgroup


The full IDs, model identities, factor-baseline comparisons, status values, source journal paths and unrounded measurements remain in the audited CSV/JSON supplement. No unfinished case enters a completed-quality comparison.


The package includes exact models, source/build definitions, frozen JAR, schedule provenance,
new raw journals and checkpoints, per-case invocations, all statuses, CSV/JSON tables,
PDF/SVG/PNG figures, paper sources and reference-verification notes. Each figure has a
case-level provenance manifest. Use a new output directory for every replay; do not overwrite
the registered outcomes. \texttt{scripts/replay\_final\_case.py} verifies input and JAR hashes
and applies the original watchdog. Maven dependencies are declared in \texttt{pom.xml};
Python plotting requires Matplotlib. Detailed commands are in the reproduction README.

All claims are bounded to these scenarios, budgets, synthetic costs and incident assumptions.
Repeated deterministic runs are not independent statistical samples. No statement of
cryptographic proof, operational mission compatibility, universal tie-rule robustness or
global mission optimum is made.
\end{document}
